
\begin{proposition}[The arithmetic genus of conductor images from coherent duality]\label{cdp:st:A.5.12}
Let $f:Y\to C_0$ be a proper connected-fibre map from a
smooth compact complex threefold with $H^2(Y,\Z)=0$ to
a smooth curve, and let its smooth general fibre have
trivial canonical bundle. This includes complex
two-tori and primary Kodaira surfaces.
Let $mS$ be a full fibre with irreducible
reduced nonnormal support $S$. Write $\eta:T\to S$ for
normalization and $\pi:Z\to T$ for a resolution.
If the pullback homomorphism
$\Pic(Y)\to\Pic(Z)$ is nonconstant, then every reduced
irreducible component $D$ of the conductor image scheme
has arithmetic genus at most one.
In particular every such component whose normalization
has positive genus is a smooth elliptic curve.
The conductor image scheme itself is allowed to be
nonreduced, and no local-complete-intersection assumption
is imposed on its reduced components.
\end{proposition}
\begin{proof}
First we determine exactly the needed canonical line on
the reduced support, without a multiplicity formula for
the direct image. The sheaf $f_*K_Y$ is torsion-free of
generic rank one on the smooth curve and hence is a
line bundle. Choose a frame near $c=f(S)$ and evaluate
it in $K_Y$. A coordinate $t$ on the base disc supplies
the conormal frame $dt$, and adjunction identifies
$K_Y|_F$ with $K_F$ after dividing by this frame.
Under that line-bundle isomorphism, the evaluated
section restricts on a general smooth fibre to a
nonzero canonical two-form, which is nowhere zero.
Its zero divisor is therefore
vertical. Shrinking the base disc to remove the other
vertical zero divisors, irreducibility of the support
gives an effective divisor $kS$, with $k\ge0$, and an
actual line-bundle comparison on $U=f^{-1}(\Delta)$:
\[
 K_Y|_U\simeq\OO_U(kS),\qquad
 \omega_S\simeq\OO_Y((k+1)S)|_S,\qquad
 \omega_S^{\otimes m}\simeq\OO_S.                 \tag{A.CD.1}\label{cdp:eq:A.CD.1}
\]
The last isomorphism uses the local base equation whose
zero divisor is exactly the full fibre $mS$.
In particular no assertion $k=m-1$ is used.

Hirzebruch--Riemann--Roch in
$0\to\OO_Y(-S)\to\OO_Y\to\OO_S\to0$ gives
$\chi(\OO_S)=0$: the rational Chern character of
$\OO_Y(-S)$ equals that of $\OO_Y$, since its first
Chern class is zero. A nonzero section of the torsion
line $\omega_S$ trivializes it. Indeed its $m$th power
is a nonzero holomorphic function on the compact
connected reduced space and hence is a nonzero constant.
Thus, with $e=h^0(S,\omega_S)$, we have
\[
 e\in\{0,1\},\qquad
 e=1\Longleftrightarrow\omega_S\simeq\OO_S,
 \qquad h^1(S,\OO_S)=1+e.                         \tag{A.CD.2}\label{cdp:eq:A.CD.2}
\]
Here $h^0(\OO_S)=1$ and Serre duality identifies
$h^2(\OO_S)$ with $e$.

Let $\mathcal Q=\eta_*\OO_T/\OO_S$ and let
$\Gamma\subset S$ be the conductor image scheme.
Both $S$ and $T$ are connected and compact, so their
global holomorphic functions are the same constants.
The normalization sequence consequently gives an
injection and the indicated kernel:
\[
 H^0(S,\mathcal Q)\hookrightarrow H^1(S,\OO_S),
 \qquad
 \operatorname{im}H^0(S,\mathcal Q)
 =\ker\bigl(H^1(S,\OO_S)\to H^1(T,\OO_T)\bigr).
                                                               \tag{A.CD.3}\label{cdp:eq:A.CD.3}
\]
Because $H^2(Y,\Z)=0$, $\Pic(Y)$ is the connected
exponential image of $H^1(Y,\OO_Y)$. A nonconstant
holomorphic group homomorphism out of it has nonzero
differential. The differential of the stated Picard
pullback is the actual composite
\[
 H^1(Y,\OO_Y)\longrightarrow H^1(S,\OO_S)
 \longrightarrow H^1(T,\OO_T)
 \longrightarrow H^1(Z,\OO_Z).
\]
Thus its middle arrow has positive rank. Equations
\eqref{cdp:eq:A.CD.2}--\eqref{cdp:eq:A.CD.3} imply $h^0(\mathcal Q)\le e$.

The conductor is Cohen--Macaulay of dimension one,
and the exact duality comparison (A.RC.3) applies
without a reducedness assumption on $\Gamma$:
$\mathcal Q=\omega_\Gamma\otimes\omega_S^{-1}|_\Gamma$.
Serre duality on that curve therefore gives
\[
 h^1(\Gamma,\omega_S|_\Gamma)
       =h^0(S,\mathcal Q)\le e.                   \tag{A.CD.4}\label{cdp:eq:A.CD.4}
\]
For any reduced irreducible component $D$ of $\Gamma$,
restriction of the line gives an exact quotient
$\omega_S|_\Gamma\to\omega_S|_D$. Its kernel is a
coherent sheaf on a one-dimensional compact space and
has no $H^2$, so the corresponding map on $H^1$ is
surjective. Hence
$h^1(D,\omega_S|_D)\le e$.

For $e=1$, \eqref{cdp:eq:A.CD.2} makes this precisely
$p_a(D)=h^1(D,\OO_D)\le1$.
For $e=0$, the same inequality makes
$h^1(D,\omega_S|_D)=0$.
The restriction $\omega_S|_D$ is torsion by \eqref{cdp:eq:A.CD.1}
and has degree zero. Curve Riemann--Roch now gives
\[
 0\le h^0(D,\omega_S|_D)
       =\chi(D,\omega_S|_D)=1-p_a(D),
\]
again proving $p_a(D)\le1$.
Finally, for a reduced irreducible curve the
normalization exact sequence gives
$p_a(D)=g(D^\nu)+\operatorname{length}
(\nu_*\OO_{D^\nu}/\OO_D)$.
Positive normalization genus therefore forces genus
one and zero normalization quotient. A normal complex
curve is smooth, which proves the final assertion.
\end{proof}


